\scanpage{91}

\uline{Proposition} Soit $f:X\unsure{\text{---}|\to}Y$ \deleted{morphisme} application constructible d'ens. schématiques. Pour tte partie constructible $E$ de $Y$, $f^{-1}(E)$ est constructible.

\diagram{\begin{tikzcd}[ampersand replacement=\&,row sep=tiny,execute at end picture={\draw[thin] (\tikzcdmatrixname-1-2.west) -- (\tikzcdmatrixname-2-1.center) -- (\tikzcdmatrixname-1-1.east); \draw[thin] (\tikzcdmatrixname-2-1.west) -- (\tikzcdmatrixname-1-1.east);}]
\deleted{X}\arrow[r,"{\deleted{\unsure{f|{}}}}"]\&\deleted{Y}\\
\deleted{X_i}\arrow[u,phantom,"{\deleted{\cup}}" description]\&
\end{tikzcd}}

En effet, soit $X=\coprod_i X_i$, les $X_i$ loc. fermés et constructibles dans $X$, \deleted{\unsure{donc|tels que}} $f_i=(f\mid |X_i|)$ un morphisme. On a $f^{-1}(E)=\bigcup f_i^{-1}(E\unsure{f|i|{}})$, il suffit de prouver que chaque $f^{-1}(E)=f^{-1}(E)\cap X_i$ est constructible dans $\alpha$, ou -- ce qui revient au $\overline{\mathrm m}$, dans $X_i$. [\deleted{\unsure{par}} \unsure{faire|par} un th. de Chevalley, \inserted{\unsure{l'immers.}} $X_i\unsure{\text{---}|\to}X$ étant de présentation finie, pour une structure induite convenable sur $X_i$ \ldots] Cela nous ramène au cas d'un \uline{morphisme} \inserted{\uline{de préschémas}} ; comme il est quasi-compact quasi-séparé, c'est connu.

% slanted
\margin{\unsure{fragile|fragilité|toujours} \unsure{pour|point} l'argument!}

\uline{Corollaire} Invariance de la notion de partie constructible d'un ens schématique

\uline{Proposition} Soit $f:X\longrightarrow Y$ une application constructible de \deleted{ens. schématiques}\inserted{préschémas quasi compacts quasi-séparés}. Pour que $f$ soit un isom d'ens schém., il f. et s. que il existe \deleted{\unsure{de|pour}} un \deleted{\unsure{p}} ens fini $I$, deux partitions $(X_i)_{i\in I}$ et $(Y_j)_{j\in I}$ de $X$ ($Y$) en parties loc. fermées constructibles, \inserted{[affines si on veut]}, telles que $f$ induise un \uline{isom} \uline{de} $X_i$ \uline{sur} $Y_i$.

C'est \unsure{évidem\textsuperscript t|évident} suffisant, \unsure{prouvons|pour} la nécessité. Soit $g$ l'inverse de $f$, $X_\alpha,Y_\beta$ des partitions \inserted{de $X,Y$} mettant en évidence le caractère constructible de $f,g$. \unsure{Évidem\textsuperscript t|Evidemment} $f,g$ sont des bijections inv. l'une de l'autre

\diagram{\begin{tikzcd}[ampersand replacement=\&]
X\arrow[r,shift left,"f"]\&Y\arrow[l,shift left,"g"]
\end{tikzcd}}
% archive sheet number: 90
